\documentclass[UTF8,10pt]{ctexart}
\usepackage[a4paper,top=2.0cm,bottom=2.0cm,left=1.55cm,right=1.55cm]{geometry}
\usepackage{fontspec,booktabs,longtable,array,xcolor,enumitem,hyperref,listings,titlesec,fancyhdr,tabularx}
\setmainfont{Times New Roman}
\setsansfont{Arial}
\setmonofont{Consolas}
\setCJKmainfont{Microsoft YaHei}
\newfontfamily\symbolfont{Segoe UI Emoji}
\definecolor{syneraBlue}{HTML}{2E74B5}
\definecolor{syneraDark}{HTML}{0B2545}
\definecolor{codeGray}{HTML}{F7F7F7}
\hypersetup{colorlinks=true,linkcolor=syneraBlue,urlcolor=syneraBlue}
\titleformat{\section}{\Large\bfseries\color{syneraBlue}}{}{0pt}{}
\titleformat{\subsection}{\large\bfseries\color{syneraDark}}{}{0pt}{}
\titleformat{\subsubsection}{\normalsize\bfseries\color{syneraDark}}{}{0pt}{}
\setlist[itemize]{leftmargin=1.8em,itemsep=0.15em,topsep=0.2em}
\setlist[enumerate]{leftmargin=1.8em,itemsep=0.15em,topsep=0.2em}
\lstset{basicstyle=\ttfamily\small,breaklines=true,frame=single,backgroundcolor=\color{codeGray},columns=fullflexible,keepspaces=true}
\pagestyle{fancy}
\fancyhf{}
\setlength{\headheight}{14pt}
\lhead{Synera: Synergy Auto-Arena}
\rhead{README}
\cfoot{\thepage}
\renewcommand{\headrulewidth}{0.4pt}
\setlength{\parindent}{0pt}
\setlength{\parskip}{0.46em}
\begin{document}
\begin{center}{\Huge\bfseries\color{syneraDark} Synera: Synergy Auto-Arena}\end{center}

\begin{quote}\itshape C++17 · Qt 6 · 单机 PVE 自走棋 · 高程课程作业\end{quote}

学号：251880336
姓名：李响

已上传github：https://github.com/winter-warm/slay\_the\_synera\_cpp\_homework
\par\medskip\hrule\medskip

\section*{目录}

\begin{itemize}
\item Part 1 —— PA 文档完成度说明
\item Part 2 —— Architecture：模块划分与文件职责
\item 2.1 总体分层
\item 2.2 程序入口与 App 流程层
\item 2.3 Core / Entity：棋盘、单位与组件
\item 2.4 Combat：战斗、Buff、Effect 与装备
\item 2.5 World：地图、事件、海克斯与背包
\item 2.6 GUI：页面、棋盘图元与面板
\item 2.7 资源与配置文件
\item 2.8 关键调用链
\item Part 3 —— 项目亮点
\item 3.1 ECS 组件模式：高解耦的角色架构
\item 3.2 灵活的 Buff / Effect 系统
\item 3.3 仿杀戮尖塔的完整玩法闭环
\item Part 4 —— AI 使用说明
\item 构建运行
\end{itemize}

\par\medskip\hrule\medskip

\section*{Part 1 —— PA 文档完成度说明}

对照 PA 说明文档中的各项需求，完成情况如下：

\subsection*{核心玩法（100\% 完成）}

{\small
\begin{longtable}{p{0.22\linewidth}p{0.20\linewidth}p{0.50\linewidth}}
\toprule
\textbf{PA 需求} & \textbf{完成状态} & \textbf{说明} \\
\midrule
M × N 六边形棋盘（8×8） & {\symbolfont ✅} & \texttt{Board} + \texttt{Hex} 实现完整六边形坐标系统与邻居/距离计算 \\
Bench 备战区（8 格） & {\symbolfont ✅} & 支持拖拽布阵、上阵上限校验 \\
Unit 属性：HP / ATK / Range / Max Mana / Mana & {\symbolfont ✅} & \texttt{StatsComponent} 完整实现，另扩展了 Defense / Shield \\
阵营区分 PlayerCtrl / EnemyCtrl & {\symbolfont ✅} & \texttt{TeamComponent} + \texttt{teams} 枚举 \\
羁绊系统（traits → 上阵阈值触发） & {\symbolfont ✅} & \texttt{Bond} 枚举含战士/法师/动物/建筑/牧师/魔女/野兽/棋灵/刺客共 9 种 \\
Prep → Combat → Resolve 三阶段 & {\symbolfont ✅} & MapPage → BattlePage（含备战）→ 自动战斗 → 结算 \\
自动战斗（索敌/移动/攻击/技能） & {\symbolfont ✅} & \texttt{AttackComponent} + \texttt{MoveComponent}，含 BFS 寻路、障碍绕行 \\
GUI 完整展示 & {\symbolfont ✅} & Qt6 全套 GUI：地图页、战斗页、事件页、开始页、HUD、卡牌面板、装备面板 \\
\bottomrule
\end{longtable}
}

\subsection*{进阶功能（全部完成）}

{\small
\begin{longtable}{p{0.22\linewidth}p{0.20\linewidth}p{0.50\linewidth}}
\toprule
\textbf{功能} & \textbf{状态} & \textbf{说明} \\
\midrule
角色合成升星（3 合 1） & {\symbolfont ✅} & \texttt{mergeOwnedCharacterCards()}，同名同星合成 \\
装备系统（穿戴 + 合成） & {\symbolfont ✅} & 基础装备 + 高级装备配方合成 \\
商店系统 & {\symbolfont ✅} & 5 刷新位 + 购买 + 刷新 + 经验/等级 \\
存档 / 读档 & {\symbolfont ✅} & \texttt{SaveManager} + 多槽位 \\
通关记录 & {\symbolfont ✅} & \texttt{RecordManager} 记录胜负及阵容 \\
角色回收（返还金币） & {\symbolfont ✅} & 备战区角色回收 \\
\bottomrule
\end{longtable}
}

\subsection*{额外实现（超出 PA 范围）}

{\small
\begin{longtable}{p{0.34\linewidth}p{0.60\linewidth}}
\toprule
\textbf{功能} & \textbf{说明} \\
\midrule
{\symbolfont 🗺}{\symbolfont ️} 路线地图 & 仿杀戮尖塔的节点地图，含普通战斗/精英/Boss/事件/休息五类节点 \\
{\symbolfont 🃏} 海克斯科技 & 每层开始可选海克斯卡，提供长期 Buff / 资源 / 路线修正 \\
{\symbolfont 📜} 事件系统 & JSON 驱动的事件链，含风险收益选择、嵌套战斗、抽奖等 \\
{\symbolfont ⚔}{\symbolfont ️} Buff / Effect 双轨系统 & 10 种触发时机 + 灵活组合，支持数据驱动 \\
{\symbolfont 🛡}{\symbolfont ️} 不可选取 / 临时阵营 & 控制类 Effect，实现魅惑、潜行等机制 \\
{\symbolfont 💀} 死亡保护 / 复活 & \texttt{deathProtection}, \texttt{deathRewind} 等 Effect \\
{\symbolfont 🔥} 燃烧 / 再生 / 属性窃取 & DOT/HOT tick 类效果 \\
{\symbolfont 📊} 战斗回放信息 & \texttt{BattleScene} 完整动画 + 状态同步 \\
\bottomrule
\end{longtable}
}

\begin{quote}\itshape \textbf{总体完成度：100\% 完成 PA 全部需求，并在此基础上增加了大量仿杀戮尖塔的扩展玩法。}\end{quote}

\par\medskip\hrule\medskip

\section*{Part 2 —— Architecture：模块划分与文件职责}

本项目按“入口/流程控制 → 世界推进 → 战斗规则 → 实体组件 → Qt 表现层 → 数据资源”的方向组织。核心思路是：**UI 不直接承担规则，规则不直接承担显示，长期状态统一进 \texttt{GameState}，战斗中的临时实体由 \texttt{battlesystem} 复制和管理**。

\subsection*{2.1 总体分层}

\begin{lstlisting}
src/main.cpp       程序入口
src/app            游戏流程、存档、窗口切换、事件和战斗回流
src/core           Hex 坐标和战术 Board
src/entity         Unit、Character、Object、Obstacle 和 ECS 组件
src/combat         战斗系统、Buff/Effect、装备、战斗配置
src/world          路线地图、事件、海克斯科技、背包/商店状态
src/gui            Qt 页面、战斗场景、HUD、卡牌面板、装备面板
assets             角色、地图、事件、装备、UI 等图片资源
\end{lstlisting}

{\small
\begin{longtable}{p{0.22\linewidth}p{0.20\linewidth}p{0.50\linewidth}}
\toprule
\textbf{层次} & \textbf{主要职责} & \textbf{上层如何使用} \\
\midrule
入口层 & 启动 Qt 应用，创建主窗口 & \texttt{main.\allowbreak{}cpp} 创建 \texttt{AppWindow} \\
App 流程层 & 管理一局游戏的状态变化、页面跳转、存档读写 & GUI 发信号给 \texttt{GameManager} \\
World 层 & 生成地图、读取事件、读取海克斯配置 & \texttt{GameManager} 进入节点时调用 \\
Combat 层 & 载入战斗、复制角色、推进自动战斗、结算胜负 & \texttt{BattlePage} 调用 \texttt{battlesystem} \\
Entity/Core 层 & 定义棋盘、坐标、单位、角色组件 & Combat 和 GUI 都依赖它们 \\
GUI 层 & 页面、场景、拖拽、面板、动画表现 & 只展示状态并把用户操作转成信号 \\
数据资源层 & JSON 配置和图片素材 & Repository / Factory 按需加载 \\
\bottomrule
\end{longtable}
}

\subsection*{2.2 程序入口与 App 流程层}

{\small
\begin{longtable}{p{0.34\linewidth}p{0.60\linewidth}}
\toprule
\textbf{文件} & \textbf{职责} \\
\midrule
\texttt{src/\allowbreak{}main.\allowbreak{}cpp} & 程序入口。创建 \texttt{QApplication}，设置应用名和版本，创建 \texttt{AppWindow} 并全屏显示。 \\
\texttt{src/\allowbreak{}app/\allowbreak{}appwindow.\allowbreak{}h} & 主窗口声明。持有 \texttt{GameManager}、\texttt{QStackedWidget}、开始页、地图页、事件页、战斗页、卡牌面板、装备面板等对象。 \\
\texttt{src/\allowbreak{}app/\allowbreak{}appwindow.\allowbreak{}cpp} & 主窗口实现。负责页面加入 \texttt{QStackedWidget}、连接 Qt 信号槽、在地图/事件/战斗/开始页之间切换，并处理居中提示 toast。 \\
\texttt{src/\allowbreak{}app/\allowbreak{}gamemanager.\allowbreak{}h} & 游戏流程中枢接口。声明新游戏、读档、进节点、选事件、选海克斯、休息、战斗结算、买卡、合成、装备等 slot。 \\
\texttt{src/\allowbreak{}app/\allowbreak{}gamemanager.\allowbreak{}cpp} & 游戏流程中枢实现。管理 \texttt{GameState}，驱动地图推进、事件动作执行、战斗启动/回流、商店刷新、自动保存和通关/失败记录。 \\
\texttt{src/\allowbreak{}app/\allowbreak{}gamestate.\allowbreak{}h} & 全局状态结构。保存 seed、玩家生命/金币、地图、已完成节点、当前事件、当前战斗、海克斯、角色卡、装备、商店等。 \\
\texttt{src/\allowbreak{}app/\allowbreak{}gamestate.\allowbreak{}cpp} & \texttt{GameState} 的 JSON 序列化/反序列化。把事件动作、战斗配置、地图、卡牌、装备、商店等状态写入/读出 JSON。 \\
\texttt{src/\allowbreak{}app/\allowbreak{}savemanager.\allowbreak{}h} & 存档管理接口。定义保存、读取、删除、查询槽位等操作。 \\
\texttt{src/\allowbreak{}app/\allowbreak{}savemanager.\allowbreak{}cpp} & 存档管理实现。把 \texttt{GameState} 写入程序目录下的 \texttt{save.\allowbreak{}json}，并从 JSON 恢复游戏。 \\
\texttt{src/\allowbreak{}app/\allowbreak{}recordmanager.\allowbreak{}h} & 战绩记录接口。定义通关/失败记录的数据结构和读写方法。 \\
\texttt{src/\allowbreak{}app/\allowbreak{}recordmanager.\allowbreak{}cpp} & 战绩记录实现。把一局游戏的结果、用时和阵容等信息记录到本地。 \\
\texttt{src/\allowbreak{}app/\allowbreak{}eventrules.\allowbreak{}h} & 事件规则辅助声明。提供事件条件、战斗上下文、奖励/惩罚等规则相关函数接口。 \\
\texttt{src/\allowbreak{}app/\allowbreak{}eventrules.\allowbreak{}cpp} & 事件规则辅助实现。处理事件条件判断、节点战斗配置、战斗失败扣血等规则细节。 \\
\bottomrule
\end{longtable}
}

App 层的核心调用关系：

\begin{lstlisting}
StartPage / MapPage / EventPage / BattlePage
        ↓ Qt signal
AppWindow
        ↓ slot 转发
GameManager
        ↓ 修改状态并 emit stateChanged
各 Page / Panel 刷新显示
\end{lstlisting}

\subsection*{2.3 Core / Entity：棋盘、单位与组件}

\subsubsection*{Core 文件}

{\small
\begin{longtable}{p{0.34\linewidth}p{0.60\linewidth}}
\toprule
\textbf{文件} & \textbf{职责} \\
\midrule
\texttt{src/\allowbreak{}core/\allowbreak{}hex.\allowbreak{}h} & 六边形坐标结构。定义 \texttt{Hex} 坐标、合法性、相等/hash 等基础能力，是棋盘寻路和距离计算的坐标基础。 \\
\texttt{src/\allowbreak{}core/\allowbreak{}board.\allowbreak{}h} & 战斗棋盘接口。定义格子、区域、占位、移动、邻居、距离、障碍物等操作。 \\
\texttt{src/\allowbreak{}core/\allowbreak{}board.\allowbreak{}cpp} & 战斗棋盘实现。读取棋盘 JSON，建立格子表，维护 \texttt{Unit*} 到 \texttt{Hex} 的映射，支持添加/移动/移除单位。 \\
\texttt{src/\allowbreak{}core/\allowbreak{}default\_\allowbreak{}board.\allowbreak{}json} & 默认战斗棋盘配置。描述棋盘格子、区域和障碍物。 \\
\bottomrule
\end{longtable}
}

\subsubsection*{Entity 基础文件}

{\small
\begin{longtable}{p{0.34\linewidth}p{0.60\linewidth}}
\toprule
\textbf{文件} & \textbf{职责} \\
\midrule
\texttt{src/\allowbreak{}entity/\allowbreak{}unit.\allowbreak{}h} & 所有棋盘单位的抽象基类。保存 id、name、position，并定义纯虚函数 \texttt{update()}。 \\
\texttt{src/\allowbreak{}entity/\allowbreak{}unit.\allowbreak{}cpp} & \texttt{Unit} 基础实现。负责单位 id 自增、位置记录等基础逻辑。 \\
\texttt{src/\allowbreak{}entity/\allowbreak{}character/\allowbreak{}character.\allowbreak{}h} & 战斗角色实体声明。组合 Stats / Attack / Move / Team / Render / Buff 等组件，并保存羁绊、稀有度、装备指针。 \\
\texttt{src/\allowbreak{}entity/\allowbreak{}character/\allowbreak{}character.\allowbreak{}cpp} & 战斗角色实体实现。定义拷贝构造、战斗开始、每帧更新顺序、死亡判断、羁绊查询等逻辑。 \\
\texttt{src/\allowbreak{}entity/\allowbreak{}character/\allowbreak{}characterfactory.\allowbreak{}h} & 角色工厂接口。按 templateId 创建角色，并提供玩家模板信息查询。 \\
\texttt{src/\allowbreak{}entity/\allowbreak{}character/\allowbreak{}characterfactory.\allowbreak{}cpp} & 角色工厂实现。从 \texttt{template.\allowbreak{}json} 读取角色模板，创建 \texttt{Character} 并装配基础属性、羁绊和默认技能。 \\
\texttt{src/\allowbreak{}entity/\allowbreak{}character/\allowbreak{}defaultskill.\allowbreak{}h} & 默认技能声明。把字符串技能名映射为 \texttt{Skill} 处理函数。 \\
\texttt{src/\allowbreak{}entity/\allowbreak{}character/\allowbreak{}defaultskill.\allowbreak{}cpp} & 默认技能实现。定义角色默认技能、主动技能和 buff 型技能的具体逻辑。 \\
\texttt{src/\allowbreak{}entity/\allowbreak{}character/\allowbreak{}template.\allowbreak{}json} & 角色模板数据。包含角色 id、名称、羁绊、稀有度、基础属性、技能名、队伍等。 \\
\texttt{src/\allowbreak{}entity/\allowbreak{}object/\allowbreak{}object.\allowbreak{}h} & 棋盘对象声明。用于非角色类棋盘单位的基础扩展。 \\
\texttt{src/\allowbreak{}entity/\allowbreak{}object/\allowbreak{}object.\allowbreak{}cpp} & 棋盘对象实现。实现对象类单位的基础行为。 \\
\texttt{src/\allowbreak{}entity/\allowbreak{}obstacle/\allowbreak{}obstacle.\allowbreak{}h} & 障碍物声明。表示棋盘上的阻挡/装饰单位。 \\
\texttt{src/\allowbreak{}entity/\allowbreak{}obstacle/\allowbreak{}obstacle.\allowbreak{}cpp} & 障碍物实现。处理障碍物的渲染资源和棋盘占位行为。 \\
\bottomrule
\end{longtable}
}

\subsubsection*{Component 文件}

{\small
\begin{longtable}{p{0.34\linewidth}p{0.60\linewidth}}
\toprule
\textbf{文件} & \textbf{职责} \\
\midrule
\texttt{src/\allowbreak{}entity/\allowbreak{}component/\allowbreak{}component.\allowbreak{}h} & 组件基类声明。保存 owner 指针和组件可用状态，是所有角色组件的公共父类。 \\
\texttt{src/\allowbreak{}entity/\allowbreak{}component/\allowbreak{}component.\allowbreak{}cpp} & 组件基类实现。提供基础构造、拷贝和启停能力。 \\
\texttt{src/\allowbreak{}entity/\allowbreak{}component/\allowbreak{}statscomponent.\allowbreak{}h} & 属性组件声明。保存 HP、MaxHP、MP、MaxMP、Attack、Defense、Shield、Range。 \\
\texttt{src/\allowbreak{}entity/\allowbreak{}component/\allowbreak{}statscomponent.\allowbreak{}cpp} & 属性组件实现。处理受击、治疗、防御减伤、护盾抵扣、死亡前 \texttt{DeathContext} 触发。 \\
\texttt{src/\allowbreak{}entity/\allowbreak{}component/\allowbreak{}attackcomponent.\allowbreak{}h} & 攻击组件声明。保存技能列表、攻击冷却、攻击间隔和行动锁。 \\
\texttt{src/\allowbreak{}entity/\allowbreak{}component/\allowbreak{}attackcomponent.\allowbreak{}cpp} & 攻击组件实现。处理普攻、攻击前后 Buff 触发、回蓝、满蓝技能释放、攻击动画请求。 \\
\texttt{src/\allowbreak{}entity/\allowbreak{}component/\allowbreak{}movecomponent.\allowbreak{}h} & 移动组件声明。定义目标选择器、寻路器、移动锁和目标缓存。 \\
\texttt{src/\allowbreak{}entity/\allowbreak{}component/\allowbreak{}movecomponent.\allowbreak{}cpp} & 移动组件实现。提供最近敌人/最高攻击/最低血量目标选择，以及 BFS 寻路和飞行寻路。 \\
\texttt{src/\allowbreak{}entity/\allowbreak{}component/\allowbreak{}teamcomponent.\allowbreak{}h} & 阵营组件声明。保存玩家/敌方等队伍信息。 \\
\texttt{src/\allowbreak{}entity/\allowbreak{}component/\allowbreak{}teamcomponent.\allowbreak{}cpp} & 阵营组件实现。支持读取、修改或临时改变队伍。 \\
\texttt{src/\allowbreak{}entity/\allowbreak{}component/\allowbreak{}rendercomponent.\allowbreak{}h} & 渲染组件声明。保存资源路径和战斗表现请求。 \\
\texttt{src/\allowbreak{}entity/\allowbreak{}component/\allowbreak{}rendercomponent.\allowbreak{}cpp} & 渲染组件实现。发出攻击突进、受击闪烁、技能爆发等表现层请求。 \\
\texttt{src/\allowbreak{}entity/\allowbreak{}component/\allowbreak{}buffcomponent.\allowbreak{}h} & Buff 组件声明。把 \texttt{BuffManager} 挂到角色身上，并暴露各种触发时机接口。 \\
\texttt{src/\allowbreak{}entity/\allowbreak{}component/\allowbreak{}buffcomponent.\allowbreak{}cpp} & Buff 组件实现。转发 onBattleStart、beforeAttack、afterAttack、beforeDeath 等事件到 \texttt{BuffManager}。 \\
\texttt{src/\allowbreak{}entity/\allowbreak{}component/\allowbreak{}interactcomponent.\allowbreak{}h} & 交互组件声明。为可交互单位预留接口。 \\
\texttt{src/\allowbreak{}entity/\allowbreak{}component/\allowbreak{}interactcomponent.\allowbreak{}cpp} & 交互组件实现。处理基础交互行为。 \\
\texttt{src/\allowbreak{}entity/\allowbreak{}component/\allowbreak{}spritecomponent.\allowbreak{}h} & 精灵组件声明。保存精灵图像资源信息。 \\
\texttt{src/\allowbreak{}entity/\allowbreak{}component/\allowbreak{}spritecomponent.\allowbreak{}cpp} & 精灵组件实现。处理精灵资源相关逻辑。 \\
\texttt{src/\allowbreak{}entity/\allowbreak{}equipskill/\allowbreak{}equipskill.\allowbreak{}h} & 装备技能接口。声明装备给角色添加开局 Buff 或主动技能的方法。 \\
\texttt{src/\allowbreak{}entity/\allowbreak{}equipskill/\allowbreak{}equipskill.\allowbreak{}cpp} & 装备技能实现。根据装备定义向角色注入 Buff、技能或属性效果。 \\
\bottomrule
\end{longtable}
}

\subsection*{2.4 Combat：战斗、Buff、Effect 与装备}

{\small
\begin{longtable}{p{0.34\linewidth}p{0.60\linewidth}}
\toprule
\textbf{文件} & \textbf{职责} \\
\midrule
\texttt{src/\allowbreak{}combat/\allowbreak{}battle.\allowbreak{}json} & 战斗配置数据。定义不同层、不同节点类型的敌人组合、棋盘 id、属性倍率等。 \\
\texttt{src/\allowbreak{}combat/\allowbreak{}battlerepository.\allowbreak{}h} & 战斗配置仓库接口。按地图节点上下文查询对应 \texttt{BattleConfig}。 \\
\texttt{src/\allowbreak{}combat/\allowbreak{}battlerepository.\allowbreak{}cpp} & 战斗配置仓库实现。从 \texttt{battle.\allowbreak{}json} 加载战斗池，按层数和战斗类型选择敌人配置。 \\
\texttt{src/\allowbreak{}combat/\allowbreak{}battlesystem.\allowbreak{}h} & 战斗系统接口。声明载入战斗、设置光环、统计羁绊、从备战单位开始战斗、逐帧更新等方法。 \\
\texttt{src/\allowbreak{}combat/\allowbreak{}battlesystem.\allowbreak{}cpp} & 战斗系统实现。复制上阵角色、生成敌人、应用装备/羁绊/海克斯、推进自动战斗、移除死亡单位、判断胜负。 \\
\texttt{src/\allowbreak{}combat/\allowbreak{}board\_\allowbreak{}utility.\allowbreak{}h} & 棋盘工具接口。提供战斗棋盘相关辅助函数。 \\
\texttt{src/\allowbreak{}combat/\allowbreak{}board\_\allowbreak{}utility.\allowbreak{}cpp} & 棋盘工具实现。处理战斗位置、棋盘辅助查询等逻辑。 \\
\texttt{src/\allowbreak{}combat/\allowbreak{}context/\allowbreak{}combatcontext.\allowbreak{}h} & 战斗上下文结构。定义 \texttt{AttackContext}、\texttt{BeAttackedContext}、\texttt{HealContext}、\texttt{DeathContext} 等 Effect 传参对象。 \\
\bottomrule
\end{longtable}
}

\subsubsection*{Buff 文件}

{\small
\begin{longtable}{p{0.34\linewidth}p{0.60\linewidth}}
\toprule
\textbf{文件} & \textbf{职责} \\
\midrule
\texttt{src/\allowbreak{}combat/\allowbreak{}buff/\allowbreak{}buff.\allowbreak{}h} & Buff 容器声明。保存 Buff 名称、持续时间、冷却、触发次数、触发时机和内部 effect 列表。 \\
\texttt{src/\allowbreak{}combat/\allowbreak{}buff/\allowbreak{}buff.\allowbreak{}cpp} & Buff 容器实现。负责 apply/remove、触发各类 effect、更新时间、处理到期和移除请求。 \\
\texttt{src/\allowbreak{}combat/\allowbreak{}buff/\allowbreak{}buffmanager.\allowbreak{}h} & Buff 管理器声明。保存角色身上的 Buff 列表，并声明各触发时机分发函数。 \\
\texttt{src/\allowbreak{}combat/\allowbreak{}buff/\allowbreak{}buffmanager.\allowbreak{}cpp} & Buff 管理器实现。添加 Buff 时先走 \texttt{BeforeAddBuff}，每帧更新 Buff，到期后移除。 \\
\texttt{src/\allowbreak{}combat/\allowbreak{}buff/\allowbreak{}bufffactory.\allowbreak{}h} & Buff 工厂接口。按 buffId 和持续时间创建 Buff。 \\
\texttt{src/\allowbreak{}combat/\allowbreak{}buff/\allowbreak{}bufffactory.\allowbreak{}cpp} & Buff 工厂实现。从 \texttt{bufflist.\allowbreak{}json} 读取 Buff 配置，并通过 \texttt{effectfactory} 装配 effect。 \\
\texttt{src/\allowbreak{}combat/\allowbreak{}buff/\allowbreak{}bufflist.\allowbreak{}json} & Buff 数据表。定义 Buff 名称、持续时间、冷却、触发时机、包含哪些 effect。 \\
\bottomrule
\end{longtable}
}

\subsubsection*{Effect 文件}

{\small
\begin{longtable}{p{0.34\linewidth}p{0.60\linewidth}}
\toprule
\textbf{文件} & \textbf{职责} \\
\midrule
\texttt{src/\allowbreak{}combat/\allowbreak{}effect/\allowbreak{}effect.\allowbreak{}h} & Effect 基类声明。用 \texttt{std::function} 保存 apply/remove/攻击前/受击前/死亡前等 handler。 \\
\texttt{src/\allowbreak{}combat/\allowbreak{}effect/\allowbreak{}effect.\allowbreak{}cpp} & Effect 基类实现。调用已注入的 handler，并保存 owner、effectId、参数等信息。 \\
\texttt{src/\allowbreak{}combat/\allowbreak{}effect/\allowbreak{}effectfactory.\allowbreak{}h} & Effect 工厂接口。按 effectId 和参数创建具体 effect 行为。 \\
\texttt{src/\allowbreak{}combat/\allowbreak{}effect/\allowbreak{}effectfactory.\allowbreak{}cpp} & Effect 工厂实现。根据 effectId 把 effect 绑定到 damage/heal/attribute/control/character 等命名空间的 handler。 \\
\texttt{src/\allowbreak{}combat/\allowbreak{}effect/\allowbreak{}damage\_\allowbreak{}effect.\allowbreak{}cpp} & 伤害类效果实现。例如燃烧、额外伤害、攻击后伤害等。 \\
\texttt{src/\allowbreak{}combat/\allowbreak{}effect/\allowbreak{}heal\_\allowbreak{}effect.\allowbreak{}cpp} & 治疗类效果实现。例如再生、伤害转治疗、治疗加成等。 \\
\texttt{src/\allowbreak{}combat/\allowbreak{}effect/\allowbreak{}attribute\_\allowbreak{}effect.\allowbreak{}cpp} & 属性类效果实现。例如加攻、减伤、免疫、护盾、死亡保护等。 \\
\texttt{src/\allowbreak{}combat/\allowbreak{}effect/\allowbreak{}control\_\allowbreak{}effect.\allowbreak{}cpp} & 控制类效果实现。例如不可选取、临时阵营变化、行动/移动限制等。 \\
\texttt{src/\allowbreak{}combat/\allowbreak{}effect/\allowbreak{}character\_\allowbreak{}effect.\allowbreak{}cpp} & 角色专属效果实现。例如特定角色位移、死亡回溯、专属机制等。 \\
\texttt{src/\allowbreak{}combat/\allowbreak{}effect/\allowbreak{}effectlist.\allowbreak{}json} & Effect 数据表。定义 effectId、名称和参数。 \\
\bottomrule
\end{longtable}
}

\subsubsection*{Equipment 文件}

{\small
\begin{longtable}{p{0.34\linewidth}p{0.60\linewidth}}
\toprule
\textbf{文件} & \textbf{职责} \\
\midrule
\texttt{src/\allowbreak{}combat/\allowbreak{}equipment/\allowbreak{}equipment.\allowbreak{}h} & 装备数据结构声明。定义基础装备、高级装备、装备组、合成配方等结构。 \\
\texttt{src/\allowbreak{}combat/\allowbreak{}equipment/\allowbreak{}equipment.\allowbreak{}cpp} & 装备基础实现。处理装备结构的辅助逻辑。 \\
\texttt{src/\allowbreak{}combat/\allowbreak{}equipment/\allowbreak{}equipmentrepository.\allowbreak{}h} & 装备仓库接口。查询基础装备、高级装备和合成结果。 \\
\texttt{src/\allowbreak{}combat/\allowbreak{}equipment/\allowbreak{}equipmentrepository.\allowbreak{}cpp} & 装备仓库实现。从 JSON 读取装备定义，支持按 group/id 查找装备和配方。 \\
\texttt{src/\allowbreak{}combat/\allowbreak{}equipment/\allowbreak{}basic\_\allowbreak{}equipment.\allowbreak{}json} & 基础装备配置。定义各基础装备属性和效果。 \\
\texttt{src/\allowbreak{}combat/\allowbreak{}equipment/\allowbreak{}advanced\_\allowbreak{}equipment.\allowbreak{}json} & 高级装备配置。定义高级装备、效果和合成配方。 \\
\bottomrule
\end{longtable}
}

\subsection*{2.5 World：地图、事件、海克斯与背包}

{\small
\begin{longtable}{p{0.34\linewidth}p{0.60\linewidth}}
\toprule
\textbf{文件} & \textbf{职责} \\
\midrule
\texttt{src/\allowbreak{}world/\allowbreak{}map/\allowbreak{}mapnode.\allowbreak{}h} & 地图节点结构。定义节点 id、层数、行列、类型、事件 id、坐标和后继节点。 \\
\texttt{src/\allowbreak{}world/\allowbreak{}map/\allowbreak{}maplayer.\allowbreak{}h} & 地图层声明。保存一层中的节点、行结构、起点/终点，并提供节点查询。 \\
\texttt{src/\allowbreak{}world/\allowbreak{}map/\allowbreak{}maplayer.\allowbreak{}cpp} & 地图层实现。负责按 id 查节点、重建 next 指针等。 \\
\texttt{src/\allowbreak{}world/\allowbreak{}map/\allowbreak{}mapgraph.\allowbreak{}h} & 地图图结构声明。保存多层地图，并提供层查询和 JSON 序列化接口。 \\
\texttt{src/\allowbreak{}world/\allowbreak{}map/\allowbreak{}mapgraph.\allowbreak{}cpp} & 地图图结构实现。处理整张地图的保存/读取和节点类型字符串转换。 \\
\texttt{src/\allowbreak{}world/\allowbreak{}map/\allowbreak{}mapgenerator.\allowbreak{}h} & 地图生成器接口。声明按 seed 生成 \texttt{MapGraph}。 \\
\texttt{src/\allowbreak{}world/\allowbreak{}map/\allowbreak{}mapgenerator.\allowbreak{}cpp} & 地图生成器实现。生成三层路线图，随机节点数和类型，固定特殊事件，并保证相邻行连通。 \\
\texttt{src/\allowbreak{}world/\allowbreak{}event/\allowbreak{}eventtypes.\allowbreak{}h} & 事件与战斗类型定义。定义 \texttt{EventActionType}、\texttt{BattleConfig}、\texttt{EventOption}、\texttt{EventDefinition} 等结构。 \\
\texttt{src/\allowbreak{}world/\allowbreak{}event/\allowbreak{}eventrepository.\allowbreak{}h} & 事件仓库接口。按事件上下文和 stepId 查询事件定义。 \\
\texttt{src/\allowbreak{}world/\allowbreak{}event/\allowbreak{}eventrepository.\allowbreak{}cpp} & 事件仓库实现。从 \texttt{events.\allowbreak{}json} 读取事件，解析步骤、选项、动作和条件，并做基础校验。 \\
\texttt{src/\allowbreak{}world/\allowbreak{}event/\allowbreak{}events.\allowbreak{}json} & 事件数据表。定义地图事件、休息事件、剧情步骤、选项和动作链。 \\
\texttt{src/\allowbreak{}world/\allowbreak{}hextech/\allowbreak{}hextech.\allowbreak{}h} & 海克斯数据结构声明。定义海克斯卡牌、效果、已选择记录和长期光环。 \\
\texttt{src/\allowbreak{}world/\allowbreak{}hextech/\allowbreak{}hextech.\allowbreak{}cpp} & 海克斯辅助实现。判断海克斯效果类型、处理记录和长期效果。 \\
\texttt{src/\allowbreak{}world/\allowbreak{}hextech/\allowbreak{}hextechrepository.\allowbreak{}h} & 海克斯仓库接口。按层数生成候选卡，按 id 查找定义。 \\
\texttt{src/\allowbreak{}world/\allowbreak{}hextech/\allowbreak{}hextechrepository.\allowbreak{}cpp} & 海克斯仓库实现。从 \texttt{hextech.\allowbreak{}json} 读取卡牌，按层数和 seed 生成 3 选 1。 \\
\texttt{src/\allowbreak{}world/\allowbreak{}hextech/\allowbreak{}hextech.\allowbreak{}json} & 海克斯配置。定义卡牌标题、描述、图片和长期/即时效果。 \\
\texttt{src/\allowbreak{}world/\allowbreak{}bag/\allowbreak{}bag.\allowbreak{}h} & 背包与商店状态结构。定义拥有的角色卡、装备实例、商店 offer、商店等级/经验等。 \\
\texttt{src/\allowbreak{}world/\allowbreak{}aura/\allowbreak{}aura.\allowbreak{}h} & 光环结构。保存战斗中由海克斯等来源提供的长期效果。 \\
\bottomrule
\end{longtable}
}

\subsection*{2.6 GUI：页面、棋盘图元与面板}

GUI 层的职责是展示和交互，不直接持久化核心规则。页面收到 \texttt{GameState} 后刷新，用户操作再通过 signal 交给 \texttt{GameManager} 或 \texttt{BattlePage}。

\subsubsection*{页面文件}

{\small
\begin{longtable}{p{0.34\linewidth}p{0.60\linewidth}}
\toprule
\textbf{文件} & \textbf{职责} \\
\midrule
\texttt{src/\allowbreak{}gui/\allowbreak{}pages/\allowbreak{}startpage.\allowbreak{}h} & 开始页声明。提供新游戏、读档等信号。 \\
\texttt{src/\allowbreak{}gui/\allowbreak{}pages/\allowbreak{}startpage.\allowbreak{}cpp} & 开始页实现。绘制开始界面，处理新游戏 seed、读取存档等操作。 \\
\texttt{src/\allowbreak{}gui/\allowbreak{}pages/\allowbreak{}mappage.\allowbreak{}h} & 地图页声明。显示路线地图、顶部状态、背包/商店/装备入口和节点选择信号。 \\
\texttt{src/\allowbreak{}gui/\allowbreak{}pages/\allowbreak{}mappage.\allowbreak{}cpp} & 地图页实现。根据 \texttt{MapGraph} 绘制节点和连线，处理节点点击、保存、返回开始页等。 \\
\texttt{src/\allowbreak{}gui/\allowbreak{}pages/\allowbreak{}eventpage.\allowbreak{}h} & 事件页声明。展示事件文本、选项、海克斯选择、休息选择、招募/装备选择等 UI。 \\
\texttt{src/\allowbreak{}gui/\allowbreak{}pages/\allowbreak{}eventpage.\allowbreak{}cpp} & 事件页实现。根据 \texttt{CurrentEventState} 切换不同选择模式，并发出对应事件信号。 \\
\texttt{src/\allowbreak{}gui/\allowbreak{}pages/\allowbreak{}battlepage.\allowbreak{}h} & 战斗页声明。管理备战、开始战斗、战斗结束、宝箱奖励、面板入口等交互。 \\
\texttt{src/\allowbreak{}gui/\allowbreak{}pages/\allowbreak{}battlepage.\allowbreak{}cpp} & 战斗页实现。连接 \texttt{BattleScene} 和 \texttt{battlesystem}，收集上阵摆放，驱动战斗计时与结果显示。 \\
\bottomrule
\end{longtable}
}

\subsubsection*{战斗场景与棋盘图元}

{\small
\begin{longtable}{p{0.34\linewidth}p{0.60\linewidth}}
\toprule
\textbf{文件} & \textbf{职责} \\
\midrule
\texttt{src/\allowbreak{}gui/\allowbreak{}battle/\allowbreak{}battlescene.\allowbreak{}h} & 战斗场景声明。管理棋盘格、单位图元、备战区、拖拽、障碍物、动画和宝箱。 \\
\texttt{src/\allowbreak{}gui/\allowbreak{}battle/\allowbreak{}battlescene.\allowbreak{}cpp} & 战斗场景实现。把 \texttt{Board} 和 \texttt{Character} 状态同步成 Qt 图元，处理拖拽布阵和战斗动画。 \\
\texttt{src/\allowbreak{}gui/\allowbreak{}board/\allowbreak{}hexlayout.\allowbreak{}h} & 六边形布局声明。负责六边形坐标和屏幕坐标之间的转换。 \\
\texttt{src/\allowbreak{}gui/\allowbreak{}board/\allowbreak{}hexlayout.\allowbreak{}cpp} & 六边形布局实现。计算 hex 中心点、顶点、多边形和命中区域。 \\
\texttt{src/\allowbreak{}gui/\allowbreak{}board/\allowbreak{}griditem.\allowbreak{}h} & 棋盘格图元声明。显示单个 hex 格子的区域、悬停和可放置状态。 \\
\texttt{src/\allowbreak{}gui/\allowbreak{}board/\allowbreak{}griditem.\allowbreak{}cpp} & 棋盘格图元实现。绘制六边形格子并响应状态变化。 \\
\texttt{src/\allowbreak{}gui/\allowbreak{}board/\allowbreak{}unititem.\allowbreak{}h} & 单位图元声明。显示战斗单位头像/棋子、血条、蓝条和动画请求。 \\
\texttt{src/\allowbreak{}gui/\allowbreak{}board/\allowbreak{}unititem.\allowbreak{}cpp} & 单位图元实现。处理单位位置、受击闪烁、攻击突进、技能爆发等表现。 \\
\texttt{src/\allowbreak{}gui/\allowbreak{}board/\allowbreak{}benchslotitem.\allowbreak{}h} & 备战区格子声明。表示一个可放置/拖拽的备战槽位。 \\
\texttt{src/\allowbreak{}gui/\allowbreak{}board/\allowbreak{}benchslotitem.\allowbreak{}cpp} & 备战区格子实现。处理备战角色显示、拖拽和选中效果。 \\
\texttt{src/\allowbreak{}gui/\allowbreak{}board/\allowbreak{}propitem.\allowbreak{}h} & 道具/场景物图元声明。用于显示障碍物、宝箱等棋盘元素。 \\
\texttt{src/\allowbreak{}gui/\allowbreak{}board/\allowbreak{}propitem.\allowbreak{}cpp} & 道具/场景物图元实现。绘制图片资源并同步位置。 \\
\bottomrule
\end{longtable}
}

\subsubsection*{HUD 与面板}

{\small
\begin{longtable}{p{0.34\linewidth}p{0.60\linewidth}}
\toprule
\textbf{文件} & \textbf{职责} \\
\midrule
\texttt{src/\allowbreak{}gui/\allowbreak{}HUD/\allowbreak{}characterhud.\allowbreak{}h} & 角色 HUD 声明。显示角色属性、技能或状态信息。 \\
\texttt{src/\allowbreak{}gui/\allowbreak{}HUD/\allowbreak{}characterhud.\allowbreak{}cpp} & 角色 HUD 实现。根据角色数据刷新 HUD。 \\
\texttt{src/\allowbreak{}gui/\allowbreak{}widgets/\allowbreak{}gamehud.\allowbreak{}h} & 游戏 HUD 声明。显示生命、金币、层数、用时等全局信息。 \\
\texttt{src/\allowbreak{}gui/\allowbreak{}widgets/\allowbreak{}gamehud.\allowbreak{}cpp} & 游戏 HUD 实现。把 \texttt{GameState} 中的核心数值同步到界面。 \\
\texttt{src/\allowbreak{}gui/\allowbreak{}widgets/\allowbreak{}cardpanel.\allowbreak{}h} & 卡牌/商店面板声明。展示拥有角色卡、商店 offer、合成/购买/训练/回收等操作。 \\
\texttt{src/\allowbreak{}gui/\allowbreak{}widgets/\allowbreak{}cardpanel.\allowbreak{}cpp} & 卡牌/商店面板实现。处理卡牌列表、商店购买、升星合成、训练选择和面板动画。 \\
\texttt{src/\allowbreak{}gui/\allowbreak{}widgets/\allowbreak{}equipmentpanel.\allowbreak{}h} & 装备面板声明。展示装备背包、穿戴、卸下、合成等操作。 \\
\texttt{src/\allowbreak{}gui/\allowbreak{}widgets/\allowbreak{}equipmentpanel.\allowbreak{}cpp} & 装备面板实现。处理装备实例、合成配方、装备到角色卡、战斗准备模式下的装备选择。 \\
\bottomrule
\end{longtable}
}

\subsection*{2.7 资源与配置文件}

{\small
\begin{longtable}{p{0.34\linewidth}p{0.60\linewidth}}
\toprule
\textbf{路径} & \textbf{职责} \\
\midrule
\texttt{assets/\allowbreak{}characters/\allowbreak{}} & 角色卡面和棋子图。角色模板中的资源路径会指向这里。 \\
\texttt{assets/\allowbreak{}events/\allowbreak{}} & 事件背景图。事件 JSON 的 \texttt{backgroundPath} 使用这些图。 \\
\texttt{assets/\allowbreak{}maps/\allowbreak{}} & 路线地图背景。每一层地图有对应背景。 \\
\texttt{assets/\allowbreak{}nodes/\allowbreak{}} & 地图节点图标。普通战斗、精英、Boss、事件、休息等节点使用不同图标。 \\
\texttt{assets/\allowbreak{}equipment/\allowbreak{}} & 装备图标。装备面板和奖励展示使用。 \\
\texttt{assets/\allowbreak{}obstacles/\allowbreak{}} & 战斗棋盘障碍物图。棋盘 JSON 会引用这些资源。 \\
\texttt{assets/\allowbreak{}pages/\allowbreak{}} & 开始页等页面背景。 \\
\texttt{assets/\allowbreak{}ui/\allowbreak{}} & 按钮、标题、HUD、宝箱等 UI 图片资源。 \\
\texttt{CMakeLists.\allowbreak{}txt} & 构建脚本。收集 \texttt{src} 下 C++ 文件，链接 Qt，并在构建后复制 JSON 和 assets 到运行目录。 \\
\bottomrule
\end{longtable}
}

运行时配置复制关系：

\begin{lstlisting}
src/core/default_board.json                 -> build/boards/default_board.json
src/world/event/events.json                 -> build/events/events.json
src/combat/battle.json                      -> build/battle/battle.json
src/combat/equipment/*.json                 -> build/equipment/*.json
src/world/hextech/hextech.json              -> build/hextech/hextech.json
assets/                                     -> build/assets/
\end{lstlisting}

\subsection*{2.8 关键调用链}

\subsubsection*{新游戏到地图}

\begin{lstlisting}
main.cpp
  -> AppWindow
  -> StartPage::newGameRequested
  -> GameManager::newGame(seed)
  -> MapGenerator::generate(seed)
  -> GameManager::refreshShopOffers()
  -> GameManager::unlockLayerStart(1)
  -> AppWindow 切换到 MapPage
\end{lstlisting}

\subsubsection*{点击地图节点}

\begin{lstlisting}
MapPage::nodeSelected(nodeId)
  -> GameManager::enterMapNode(nodeId)
  -> 根据 event_id 分流：
       Start        -> loadHexTechEvent()
       Rest         -> loadRestEvent()
       Normal/Elite -> BattleRepository::battleFor() -> startBattle()
       Boss         -> BattleRepository::battleFor() -> startBattle()
       Event        -> EventRepository::eventFor() -> EventPage
\end{lstlisting}

\subsubsection*{一场战斗}

\begin{lstlisting}
BattlePage 收集备战区 placements
  -> GameManager::startPreparedBattle(placements)
  -> battlesystem::startFromPreparedUnits()
       复制玩家角色和敌人
       应用装备 Buff
       应用羁绊加成
       应用海克斯光环
       放入 Board
  -> battlesystem::update(deltaTime)
       Character::update()
         BuffComponent::update()
         MoveComponent::update()
         AttackComponent::update()
       removeDeadCharacters()
       finishIfTeamDead()
  -> BattlePage::battleFinished
  -> GameManager::finishBattle()
\end{lstlisting}

\subsubsection*{事件动作}

\begin{lstlisting}
EventPage::eventOptionSelected(index)
  -> GameManager::chooseEventOption(index)
  -> GameManager::executeEventActions()
  -> GameManager::executeEventAction()
       FinishNode / GoToStep / StartBattle / GrantGold / LoseHp / ChooseRecruit ...
\end{lstlisting}

\par\medskip\hrule\medskip

\section*{Part 3 —— 项目亮点}

\subsection*{3.1 ECS 组件模式：高解耦的角色架构}

传统做法通常把所有属性塞进 \texttt{Character} 类，导致类膨胀且难以扩展。本项目采用了 \textbf{ECS（Entity-Component）组件化设计}：

\begin{lstlisting}
Character (Entity)
  ├── StatsComponent      —— 属性：HP/ATK/MP/Defense/Shield/Range
  ├── AttackComponent     —— 战斗：普攻冷却/技能列表/行动锁
  ├── MoveComponent       —— 移动：寻路/障碍绕行/路径缓存
  ├── TeamComponent       —— 阵营：玩家/敌方/临时转换
  ├── RenderComponent     —— 渲染：棋子图标/卡面/动画帧
  ├── BuffComponent       —— Buff：挂载/触发/到期移除
  └── Equipment*          —— 装备槽（最多 1 件）
\end{lstlisting}

\textbf{设计优势：}

\begin{itemize}
\item \textbf{单一职责}：每个 Component 只管理一个维度，互不干扰。改移动逻辑不需要看战斗代码。
\item \textbf{可拔插}：新角色只需要组合不同 Component，不需要修改基类。\texttt{Obstacle}（障碍物）复用 \texttt{Unit} 体系，但禁用了不需要的组件。
\item \textbf{可扩展}：添加新系统（如飞行、隐身）只需新增一个 Component 或在现有组件上加一个 lock 计数器，不影响已有逻辑。
\item \textbf{复制安全}：每个 Component 都提供了拷贝构造函数（\texttt{Component(const Component\&, Character*)}），合成升星时能正确复制所有状态。
\end{itemize}

\subsection*{3.2 灵活的 Buff / Effect 系统}

这是整个战斗系统的心脏，设计核心是 \textbf{Buff 作为容器 + Effect 作为行为} 的双层解耦：

\subsubsection*{Buff 层（容器）}

\begin{lstlisting}
class buff {
    QString name;
    vector<unique_ptr<effect>> effects;  // 一个 Buff 可包含多个 Effect
    float duration;           // 持续时间（-1 = 永久）
    float cooldown;           // 冷却时间
    BuffTrigger trigger;      // 触发时机
    int maxtrigger;           // 最大触发次数
    short priority;           // 优先级
};
\end{lstlisting}

\textbf{10 种触发时机}覆盖完整战斗生命周期：

{\small
\begin{longtable}{p{0.34\linewidth}p{0.60\linewidth}}
\toprule
\textbf{触发时机} & \textbf{典型用途} \\
\midrule
\texttt{OnBattleStart} & 战斗开始加属性、飞行、视野 \\
\texttt{BeforeAttack} & 攻击增伤、暴击 \\
\texttt{AfterAttack} & 攻击后附加伤害、女巫化 \\
\texttt{BeforeBeAttacked} & 伤害减免、闪避、反伤 \\
\texttt{AfterBeAttacked} & 受伤后加攻、位移、模仿 \\
\texttt{BeforeHeal} & 治疗加成 \\
\texttt{AfterSkill} & 技能后效果 \\
\texttt{BeforeAddBuff} & Buff 抗性/免疫 \\
\texttt{BeforeDeath} & 死亡保护、复活 \\
\texttt{OnTurnStart} & DOT 燃烧、HOT 再生 \\
\bottomrule
\end{longtable}
}

\subsubsection*{Effect 层（行为）}

\begin{lstlisting}
class effect {
    // 每种触发时机对应一个 std::function 回调
    AttackHandler beforeAttackFunc;      // 可通过 setBeforeAttack() 注入
    BeAttackedHandler beforeBeAttackedFunc;
    TurnHandler onTurnStartFunc;
    // ... 共 10 种 handler
};
\end{lstlisting}

\textbf{关键设计巧思：}

\begin{enumerate}
\item \textbf{函数注入而非继承} —— Effect 不使用虚函数多态，而是通过 \texttt{std::function} 注入行为。\texttt{effectfactory} 从 JSON 读取配置后，按 effectId 装配对应的 handler。这意味着添加新效果只需在 factory 的 switch 里加一个 case，不需要新建子类。
\end{enumerate}

\begin{enumerate}
\item \textbf{namespace 分组管理} —— 20+ 种 Effect 按类型分到 5 个 namespace：
\end{enumerate}
\begin{itemize}
\item \texttt{damage\_\allowbreak{}effect} —— 女巫化、燃烧伤害
\item \texttt{heal\_\allowbreak{}effect} —— 反转伤害为治疗、再生
\item \texttt{attribute\_\allowbreak{}effect} —— 属性修改、伤害倍率、减伤、免疫、死亡保护
\item \texttt{control\_\allowbreak{}effect} —— 临时阵营（魅惑）、不可选取（潜行）
\item \texttt{character\_\allowbreak{}effect} —— 角色专属：Coco 位移、Noah 减伤再生、Levitation 飞行、死亡回溯等
\end{itemize}

\begin{enumerate}
\item \textbf{Context 传参模式} —— 每个触发时机有独立 Context 结构体（\texttt{AttackContext}, \texttt{BeAttackedContext}, \texttt{DeathContext} 等），Effect 可以修改其中的 \texttt{cancelled}、\texttt{damage} 等字段来拦截/修改后续行为。例如 \texttt{damageImmunity} 通过设置 \texttt{cancelled = true} 完全阻止伤害。
\end{enumerate}

\begin{enumerate}
\item \textbf{JSON 数据驱动} —— Buff 列表 (\texttt{bufflist.\allowbreak{}json}) 和 Effect 列表 (\texttt{effectlist.\allowbreak{}json}) 完全由数据定义，策划可以直接修改 JSON 来调整效果参数，无需重新编译。
\end{enumerate}

\begin{enumerate}
\item \textbf{BuffManager 优先级排序} —— 多个 Buff 按优先级排序触发，确保关键效果（如死亡保护）在其他效果之前执行。
\end{enumerate}

\subsection*{3.3 仿杀戮尖塔的完整玩法闭环}

在 PA 要求的基础自走棋之上，补充了类杀戮尖塔的完整 Roguelike 爬塔体验：

\begin{lstlisting}
开始 → 选择海克斯 → 推进地图节点 → 战斗/事件/休息 → Boss 层 → 下一层 → 通关/失败
\end{lstlisting}

{\small
\begin{longtable}{p{0.34\linewidth}p{0.60\linewidth}}
\toprule
\textbf{系统} & \textbf{实现细节} \\
\midrule
\textbf{路线地图} & \texttt{MapGraph} 生成分层节点图，玩家逐节点推进，选择路线 \\
\textbf{海克斯科技} & 每层开始 3 选 1，效果涵盖长期 Buff、资源奖励、路线/掉落修正 \\
\textbf{事件系统} & JSON 驱动的事件链，支持多选项、条件判定、嵌套战斗、抽奖（\texttt{LotteryScratch}）等 18 种 Action \\
\textbf{休息节点} & 休息回血 / 训练升星 / 挖宝得装备 / 锻造合成，4 种选择 \\
\textbf{Boss 战} & 三层 Boss，击败即通关 \\
\textbf{存档系统} & 手动 + 自动存档，多槽位 \\
\textbf{通关记录} & 记录胜负、阵容、用时 \\
\bottomrule
\end{longtable}
}

\par\medskip\hrule\medskip

\section*{Part 4 —— AI 使用说明}

本项目在开发过程中使用了 AI（Claude Code和codex）辅助，以下是各模块的 AI 参与程度说明：

\subsection*{4.1 Entity 部分 —— 主要自己完成，AI 修 Bug和机械化劳动}

{\small
\begin{longtable}{p{0.34\linewidth}p{0.60\linewidth}}
\toprule
\textbf{模块} & \textbf{文件} \\
\midrule
Unit 基类 & \texttt{entity/\allowbreak{}unit.\allowbreak{}h/\allowbreak{}cpp} \\
Character & \texttt{entity/\allowbreak{}character/\allowbreak{}character.\allowbreak{}h/\allowbreak{}cpp} \\
各 Component & \texttt{entity/\allowbreak{}component/\allowbreak{}*.\allowbreak{}h/\allowbreak{}cpp} \\
CharacterFactory & \texttt{entity/\allowbreak{}character/\allowbreak{}characterfactory.\allowbreak{}h/\allowbreak{}cpp} \\
DefaultSkill & \texttt{entity/\allowbreak{}character/\allowbreak{}defaultskill.\allowbreak{}h/\allowbreak{}cpp} \\
Object / Obstacle & \texttt{entity/\allowbreak{}object/\allowbreak{}}, \texttt{entity/\allowbreak{}obstacle/\allowbreak{}} \\
EquipSkill & \texttt{entity/\allowbreak{}equipskill/\allowbreak{}} \\
\bottomrule
\end{longtable}
}


\subsection*{4.2 图片素材 —— 完全 AI 生成}

{\small
\begin{longtable}{p{0.22\linewidth}p{0.20\linewidth}p{0.50\linewidth}}
\toprule
\textbf{资源类型} & \textbf{路径} & \textbf{生成方式} \\
\midrule
角色卡面 + 棋子图 & \texttt{assets/\allowbreak{}characters/\allowbreak{}*/\allowbreak{}card.\allowbreak{}png}, \texttt{pawn.\allowbreak{}png} & AI 图像生成 \\
事件背景图 & \texttt{assets/\allowbreak{}events/\allowbreak{}} & AI 图像生成 \\
地图背景 & \texttt{assets/\allowbreak{}maps/\allowbreak{}} & AI 图像生成 \\
节点图标 & \texttt{assets/\allowbreak{}nodes/\allowbreak{}} & AI 图像生成 \\
装备图标 & \texttt{assets/\allowbreak{}equipment/\allowbreak{}} & AI 图像生成 \\
障碍物图 & \texttt{assets/\allowbreak{}obstacles/\allowbreak{}} & AI 图像生成 \\
UI 元素 & \texttt{assets/\allowbreak{}ui/\allowbreak{}}, \texttt{assets/\allowbreak{}pages/\allowbreak{}} & AI 图像生成 \\
\bottomrule
\end{longtable}
}

\begin{quote}\itshape 提示词在github说放了\end{quote}

\subsection*{4.3 GUI 部分 —— 大部分 AI 生成}

{\small
\begin{longtable}{p{0.20\linewidth}p{0.24\linewidth}p{0.18\linewidth}p{0.28\linewidth}}
\toprule
\textbf{模块} & \textbf{文件} & \textbf{AI 参与度} & \textbf{说明} \\
\midrule
StartPage & \texttt{gui/\allowbreak{}pages/\allowbreak{}startpage.\allowbreak{}h/\allowbreak{}cpp} & AI 生成，不太懂 Qt &  \\
MapPage & \texttt{gui/\allowbreak{}pages/\allowbreak{}mappage.\allowbreak{}h/\allowbreak{}cpp} & AI 生成 &  \\
BattlePage & \texttt{gui/\allowbreak{}pages/\allowbreak{}battlepage.\allowbreak{}h/\allowbreak{}cpp} & 在原有上改的 &  \\
EventPage & \texttt{gui/\allowbreak{}pages/\allowbreak{}eventpage.\allowbreak{}h/\allowbreak{}cpp} & 自己写的 &  \\
BattleScene & \texttt{gui/\allowbreak{}battle/\allowbreak{}battlescene.\allowbreak{}h/\allowbreak{}cpp} & 战斗棋盘 Qt 场景，AI 生成 &  \\
Board UI & \texttt{gui/\allowbreak{}board/\allowbreak{}hexlayout.\allowbreak{}h/\allowbreak{}cpp}, \texttt{unititem}, \texttt{griditem}, \texttt{benchslotitem}, \texttt{propitem} & 六边形的坐标表示是在b站学的，没有用原来的坐标 &  \\
HUD / Widgets & \texttt{gui/\allowbreak{}HUD/\allowbreak{}}, \texttt{gui/\allowbreak{}widgets/\allowbreak{}} & CardPanel, EquipmentPanel, GameHUD 等 &  \\
\bottomrule
\end{longtable}
}

\begin{quote}\itshape 由于对 Qt 的 Graphics View / Widget 体系不太熟悉，GUI 层绝大部分代码由 AI 生成，自己主要负责接口定义和决定采用的数学方法，并给出优化方案\end{quote}

\subsection*{4.4 Combat 部分 —— AI 辅助，自己搭框架}

{\small
\begin{longtable}{p{0.22\linewidth}p{0.20\linewidth}p{0.50\linewidth}}
\toprule
\textbf{模块} & \textbf{自己完成} & \textbf{AI 完成} \\
\midrule
\textbf{BattleSystem} & \texttt{.\allowbreak{}h} 头文件（接口设计 + 架构） & \texttt{.\allowbreak{}cpp} 实现文件 \\
\textbf{Effect 基类} & 设计 Effect 抽象基类 + 10 种 handler 函数指针注入模式 & — \\
\textbf{Effect 派生逻辑} & 搭建 namespace 分组 + factory 框架 & 各具体 Effect handler 实现（20+ 种） \\
\textbf{Buff} & 设计 Buff 类 + BuffManager 容器 + 10 种触发时机枚举 & — \\
\textbf{BattleRepository} & 接口设计 & \texttt{.\allowbreak{}cpp} JSON 加载实现 \\
\textbf{Board Utility} & 自己完成 & — \\
\textbf{Equipment} & 接口设计 & \texttt{.\allowbreak{}cpp} JSON 加载 \\
\textbf{Context} & \texttt{combatcontext.\allowbreak{}h} 全部 Context 结构体设计 & — \\
\bottomrule
\end{longtable}
}

\begin{quote}\itshape Combat 模块的核心架构（Effect 基类、Buff 容器、Context 传参、Factory 装配）全部自己设计。由于 Effect 种类特别多（20+ 种），具体的 handler 实现交给 AI 批量生成；BattleSystem 采用"自己写 \texttt{.\allowbreak{}h} 定义接口，AI 写 \texttt{.\allowbreak{}cpp} 实现细节"的协作模式。\end{quote}

\subsection*{4.5 其余部分}

{\small
\begin{longtable}{p{0.22\linewidth}p{0.20\linewidth}p{0.50\linewidth}}
\toprule
\textbf{模块} & \textbf{AI 参与度} & \textbf{说明} \\
\midrule
App 层 (\texttt{gamemanager}, \texttt{gamestate}, \texttt{savemanager}, \texttt{recordmanager}, \texttt{eventrules}) & 游戏流程控制基本由 AI 完成，这个太底层了我做不来，没法像ai一样记住那么多模块的关系 &  \\
World 层 (\texttt{mapgenerator}, \texttt{mapgraph}, \texttt{maplayer}, \texttt{eventrepository}, \texttt{hextechrepository}) & ai太笨了，看不懂需要绘图的map(主要是它没办法运行)最后变成了我写 &  \\
Core 层 (\texttt{board}, \texttt{hex}) & 原有，修改了一下 &  \\
\bottomrule
\end{longtable}
}

\section*{资源结构}

\begin{lstlisting}
assets/characters   角色卡面和棋子图
assets/events       事件背景图
assets/maps         路线地图背景
assets/nodes        地图节点图标
assets/equipment    装备图标
assets/obstacles    棋盘障碍物图
assets/pages        开始页背景
assets/ui           按钮、标题、HUD、宝箱等 UI 图
\end{lstlisting}

\par\medskip\hrule\medskip

\section*{构建运行}

\subsection*{环境要求}

\begin{itemize}
\item \textbf{Qt 6}（\texttt{Core}、\texttt{Gui}、\texttt{Widgets}、\texttt{Svg}、\texttt{Network} 模块）
\item \textbf{CMake} 3.16+
\item \textbf{编译器}：支持 C++17（GCC 13+ / MSVC 2022+ / MinGW-w64）
\end{itemize}

\subsection*{构建}

\begin{lstlisting}
cmake -S . -B build -DCMAKE_BUILD_TYPE=Release -DCMAKE_PREFIX_PATH="/path/to/Qt/6.x.x/mingw_64"
cmake --build build --config Release
\end{lstlisting}
\end{document}
